% ============================================================
% Thesis: Tracking the Lifetime of Malicious Indicators
%         in Public Threat Intelligence Feeds
% Author: Dimitris Keramidas
% ============================================================
\documentclass[12pt, a4paper, twoside]{report}

% --- Encoding & Language ---
\usepackage[utf8]{inputenc}
\usepackage[T1]{fontenc}
\usepackage[english, greek]{babel}
\usepackage{alphabeta}          % Allow Greek letters in math mode

% --- Typography ---
\usepackage{lmodern}
\usepackage{microtype}
\usepackage{setspace}
\onehalfspacing

% --- Page Geometry ---
\usepackage[
  a4paper,
  top=2.5cm, bottom=2.5cm,
  inner=3cm, outer=2cm,
  headheight=15pt
]{geometry}

% --- Headers / Footers ---
\usepackage{fancyhdr}
\pagestyle{fancy}
\fancyhf{}
\fancyhead[LE,RO]{\thepage}
\fancyhead[LO]{\nouppercase{\rightmark}}
\fancyhead[RE]{\nouppercase{\leftmark}}
\renewcommand{\headrulewidth}{0.4pt}

% --- Graphics & Figures ---
\usepackage{graphicx}
\usepackage{float}
\usepackage{subcaption}
\graphicspath{{../figures/}{./figures/}}

% --- Tables ---
\usepackage{booktabs}
\usepackage{tabularx}
\usepackage{longtable}
\usepackage{multirow}

% --- Mathematics ---
\usepackage{amsmath, amssymb}

% --- Code Listings ---
\usepackage{listings}
\usepackage{xcolor}
\lstset{
  basicstyle=\ttfamily\small,
  showstringspaces=false,
  breaklines=true,
  frame=single,
  numbers=left,
  numberstyle=\tiny\color{gray},
  keywordstyle=\color{blue},
  commentstyle=\color{green!50!black},
  stringstyle=\color{red!70!black},
  backgroundcolor=\color{gray!5},
}

% --- URLs & Hyperlinks ---
\usepackage[
  colorlinks=true,
  linkcolor=black,
  citecolor=blue,
  urlcolor=blue,
  bookmarks=true,
  pdfauthor={Dimitris Keramidas},
  pdftitle={Tracking the Lifetime of Malicious Indicators in Public Threat Intelligence Feeds}
]{hyperref}
\usepackage{url}

% --- Bibliography ---
\usepackage[numbers, sort&compress]{natbib}
\bibliographystyle{plainnat}

% --- Miscellaneous ---
\usepackage{enumitem}
\usepackage{csquotes}
\usepackage{pdfpages}
\usepackage[acronym, toc]{glossaries}
\makeglossaries

\newacronym{tls}{TLS}{Transport Layer Security}
\newacronym{https}{HTTPS}{Hypertext Transfer Protocol Secure}
\newacronym{c2}{C2}{command-and-control}
\newacronym{ioc}{IOC}{indicator of compromise}
\newacronym{jarm}{JARM}{a TLS server fingerprinting method by Salesforce Engineering}
\newacronym{dns}{DNS}{Domain Name System}
\newacronym{ttl}{TTL}{time to live}
\newacronym{asn}{ASN}{autonomous system number}
\newacronym{bgp}{BGP}{Border Gateway Protocol}
\newacronym{ml}{ML}{machine learning}
\newacronym{gnn}{GNN}{graph neural network}
\newacronym{api}{API}{application programming interface}
\newacronym{ip}{IP}{Internet Protocol}
\newacronym{ipv4}{IPv4}{Internet Protocol version 4}
\newacronym{ipv6}{IPv6}{Internet Protocol version 6}
\newacronym{cli}{CLI}{command-line interface}
\newacronym{csv}{CSV}{comma-separated values}
\newacronym{rq}{RQ}{research question}

% ============================================================
\begin{document}
\selectlanguage{english}

% ============================================================
%  FRONT MATTER
% ============================================================
\begin{titlepage}
  \centering

  {\Large\bfseries Technical University of Crete\\[0.3em]
  School of Electrical and Computer Engineering}\\[0.5em]
  {\large Diploma Thesis}\\[0.8cm]

  \rule{\linewidth}{0.5mm}\\[0.4cm]
  {\LARGE\bfseries Tracking the Lifetime of Malicious Indicators\\[0.3em]
  in Public Threat Intelligence Feeds}\\[0.4cm]
  \rule{\linewidth}{0.5mm}\\[0.5cm]

  {\large\itshape
  Παρακολούθηση της διάρκειας ζωής κακόβουλων δεικτών\\[0.2em]
  σε ροές πληροφοριών για δημόσιες απειλές}\\[0.9cm]

  \begin{tabular}{@{}l@{\hspace{3cm}}l@{}}
    \textit{Author:} & \textit{Thesis Committee:} \\
    Dimitris Keramidas & [Supervisor Name] (Supervisor) \\
     & [Committee Member 2] \\
     & [Committee Member 3] \\
  \end{tabular}\\[0.8cm]

  \includegraphics[width=2cm]{tuc_logo.png}\\[0.8cm]

  {\itshape A thesis submitted in fulfillment of the requirements\\
  for the diploma of Electrical and Computer Engineer\\[0.3em]
  in the\\[0.3em]
  School of Electrical and Computer Engineering}\\[0.6cm]

  \vfill
  {\large Chania, Crete, Greece \\ \the\year}
\end{titlepage}

\cleardoublepage

% --- Copyright / Declaration ---
\thispagestyle{empty}
\vspace*{\fill}
\begin{center}
  \textcopyright\ \the\year\ Dimitris Keramidas. All rights reserved.
\end{center}
\vspace*{\fill}
\cleardoublepage

% --- Abstract (English) ---
\chapter*{Abstract}
\addcontentsline{toc}{chapter}{Abstract}

Public threat intelligence feeds are a primary resource for identifying malicious activity on the
internet. This study investigates the lifetime of malicious indicators across six public feeds
--- VirusTotal, AbuseIPDB, Censys, ThreatFox, Cinsscore, and OpenPhish --- using IP addresses
classified as malicious by TLScope as ground truth. TLScope is an active TLS fingerprinting system
that identifies malicious servers through machine learning without inspecting payload content.

Each IP is tracked on a predominantly three-day cadence across all providers for one to two
months, measuring
\textbf{detection lag} (how quickly providers identify what TLScope has already found),
\textbf{indicator persistence} (how long a detection is maintained), and
\textbf{inter-provider agreement}. Across eight collection waves covering 18\,428 TLScope-identified
IPs, 81.2\% were never flagged by any of the six tracked providers, and among those that were,
agreement was low, with 83.9\% of detected IPs flagged by only a single detecting engine. These
results quantify a substantial gap between specialised active scanning and public feed coverage,
contributing to a broader understanding of threat intelligence ecosystem effectiveness and the
temporal lifecycle of malicious indicators.

\bigskip\noindent
\textbf{Keywords:} threat intelligence, indicator of compromise, malicious IP, VirusTotal,
TLS fingerprinting, detection lag, indicator lifetime, botnet

\cleardoublepage

% --- Abstract (Greek) ---
\selectlanguage{greek}
\chapter*{Περίληψη}
\addcontentsline{toc}{chapter}{Περίληψη}

Οι δημόσιες ροές πληροφοριών για απειλές αποτελούν σημαντικό εργαλείο για τον εντοπισμό
κακόβουλης δραστηριότητας στο διαδίκτυο. Στην παρούσα εργασία διερευνάται η διάρκεια ζωής
κακόβουλων δεικτών σε έξι δημόσιες ροές (\foreignlanguage{english}{VirusTotal},
\foreignlanguage{english}{AbuseIPDB}, \foreignlanguage{english}{Censys},
\foreignlanguage{english}{ThreatFox}, \foreignlanguage{english}{Cinsscore}
και \foreignlanguage{english}{OpenPhish}), χρησιμοποιώντας ως δεδομένα αναφοράς διευθύνσεις
\foreignlanguage{english}{IP} που έχουν ταξινομηθεί ως
κακόβουλες από το εργαλείο \foreignlanguage{english}{TLScope}. Κάθε διεύθυνση παρακολουθείται κατά κύριο λόγο ανά τρεις
ημέρες σε κάθε
πάροχο για διάστημα ενός έως δύο μηνών, με σκοπό τη μέτρηση της καθυστέρησης ανίχνευσης,
της επιμονής των δεικτών και του βαθμού συμφωνίας μεταξύ των παρόχων. Σε οκτώ κύκλους συλλογής
δεδομένων, που κάλυψαν 18.428 διευθύνσεις \foreignlanguage{english}{IP} ταξινομημένες από το
\foreignlanguage{english}{TLScope}, το 81,2\% δεν
επισημάνθηκε ποτέ από κανέναν από τους έξι παρόχους που παρακολουθήθηκαν, ενώ μεταξύ των
υπολοίπων η συμφωνία ήταν χαμηλή, καθώς το 83,9\% των δεικτών που εντοπίστηκαν επισημάνθηκαν από
μόνο μία μηχανή ανίχνευσης. Τα αποτελέσματα αυτά ποσοτικοποιούν ένα σημαντικό χάσμα μεταξύ
εξειδικευμένης ενεργής σάρωσης \foreignlanguage{english}{TLS} και της κάλυψης από δημόσιες ροές πληροφοριών, συμβάλλοντας
στην κατανόηση της αποτελεσματικότητας του οικοσυστήματος
\foreignlanguage{english}{Threat Intelligence}.

\bigskip\noindent
\textbf{Λέξεις-κλειδιά:} πληροφορίες απειλών, δείκτης παραβίασης, κακόβουλη
\foreignlanguage{english}{IP},
\foreignlanguage{english}{VirusTotal}, αποτύπωμα \foreignlanguage{english}{TLS}, καθυστέρηση ανίχνευσης, διάρκεια ζωής δείκτη,
\foreignlanguage{english}{botnet}

\selectlanguage{english}
\cleardoublepage

% --- Acknowledgements ---
\chapter*{Acknowledgements}
\addcontentsline{toc}{chapter}{Acknowledgements}

I would like to thank my supervisor for their guidance throughout this project,
and the authors of TLScope for providing the classified IP datasets that form the
empirical foundation of this work.

\cleardoublepage

% --- Table of Contents ---
\tableofcontents
\cleardoublepage

\listoffigures
\addcontentsline{toc}{chapter}{List of Figures}
\cleardoublepage

\listoftables
\addcontentsline{toc}{chapter}{List of Tables}
\cleardoublepage

\glsaddall
\printglossary[type=\acronymtype, title={List of Abbreviations}, nonumberlist]
\cleardoublepage

% ============================================================
%  MAIN MATTER
% ============================================================

% ============================================================
\chapter{Introduction}
\label{ch:introduction}
% ============================================================

\section{Motivation}

The modern internet threat landscape relies heavily on \emph{public threat intelligence feeds} ---
continuously updated lists of known-malicious IP addresses, domains, and URLs. Security operations
centres, firewalls, intrusion detection systems, and endpoint products consume these feeds to block
or flag traffic associated with malware, botnet command-and-control (C2) infrastructure, phishing
campaigns, and abuse activity. The practical value of a feed depends critically on two properties:
(a) its \textbf{coverage} --- the fraction of genuinely malicious indicators it detects --- and
(b) its \textbf{timeliness} --- how quickly it detects and how long it retains an indicator once
detected.

Despite the central role these feeds play in operational security, their temporal behaviour has
received limited systematic study. It is not well understood how long after a malicious server is
deployed it takes for public feeds to flag its IP address, how long the flag persists once added,
or to what extent different providers agree with each other.

\section{Problem Statement}

This thesis addresses the following research questions:

\begin{enumerate}[label=\textbf{RQ\arabic*.}]
  \item \textbf{Detection lag:} How many days pass between a server being identified as malicious
    by TLScope and its IP address being flagged by each public threat intelligence provider?
  \item \textbf{Indicator persistence:} Once a provider flags an IP, how long does it maintain
    that classification before the indicator expires?
  \item \textbf{Inter-provider agreement:} Do different providers detect the same IPs, or does
    each cover a largely distinct subset of the malicious population?
\end{enumerate}

\section{Approach}

The study uses IP addresses classified as malicious by \textbf{TLScope}
\cite{theofanous2024tlscope}, an active TLS fingerprinting system, as ground truth.
Each group of 1\,596--2\,728 malicious IPs (18\,428 IPs across eight completed groups) is tracked
on a predominantly three-day cadence (\S\ref{sec:collection_schedule}) across six providers for one
to two months. Six providers are evaluated:

\begin{itemize}
  \item \textbf{VirusTotal} --- multi-engine antivirus aggregate
  \item \textbf{AbuseIPDB} --- community abuse score
  \item \textbf{Censys} --- internet-wide scanner with C2 labels
  \item \textbf{ThreatFox} --- malware IOC feed (abuse.ch)
  \item \textbf{Cinsscore} --- CINS Army List (daily blacklist)
  \item \textbf{OpenPhish} --- phishing URL feed with DNS resolution
\end{itemize}

This is the design each wave was set up to follow; in practice, Censys tracking was discontinued
partway through the study (\S\ref{sec:provider_availability}), which the results in
Chapter~\ref{ch:results} account for explicitly.

\section{Contributions}

The main contributions of this thesis are:
\begin{enumerate}
  \item A longitudinal dataset of 18\,428 TLScope-identified malicious IPs tracked across six
    public threat intelligence providers over eight waves spanning eleven months, with a
    one-to-two-month observation window per wave.
  \item Empirical measurements of detection lag, indicator lifetime, and inter-provider agreement,
    reported both per provider and per individual detection engine.
  \item A quantified coverage gap between active TLS scanning and public feed coverage: 81.2\% of
    TLScope-identified IPs are never flagged by any tracked provider, and that rate rises
    significantly over the study period rather than holding steady.
  \item Evidence that low inter-provider agreement reflects genuine specialisation rather than
    sparse detection, obtained by comparing observed single-engine detection against a Poisson
    binomial chance baseline derived from each engine's own detection rate
    (\S\ref{sec:agreement_vs_chance}).
  \item A provider-level accounting showing that coverage is concentrated almost entirely in
    VirusTotal's aggregated engines, with one tracked provider (Cinsscore) detecting none of the
    18\,428 IPs at any point (\S\ref{sec:provider_coverage}).
  \item Open-source tooling (\texttt{NetScan}) for reproducible longitudinal IP tracking across
    multiple providers.
\end{enumerate}

\section{Thesis Structure}

The remainder of this thesis is organised as follows. Chapter~\ref{ch:background} provides
background on TLS fingerprinting and threat intelligence ecosystems.
Chapter~\ref{ch:related_work} surveys related work.
Chapter~\ref{ch:methodology} describes the data collection and analysis methodology.
Chapter~\ref{ch:providers} details each threat intelligence provider.
Chapter~\ref{ch:implementation} describes the implementation of the NetScan system.
Chapter~\ref{ch:results} presents experimental results.
Chapter~\ref{ch:discussion} interprets the findings and sets out the threats to validity that bound
them. Chapter~\ref{ch:future_work} outlines future directions.
Chapter~\ref{ch:conclusion} concludes the thesis.

% ============================================================
\chapter{Background}
\label{ch:background}
% ============================================================

\section{TLS and Encrypted Malicious Infrastructure}

The Transport Layer Security (TLS) protocol is the cryptographic foundation of HTTPS and many
other application protocols. By 2022, more than 85\% of malware used encrypted communication
channels, making payload-based detection infeasible for classifying traffic without decrypting
it~\cite{theofanous2024tlscope}. TLS metadata --- cipher suite selection, TLS version, certificate
properties --- remains visible even when payload content is encrypted, providing an alternative
detection surface.

\subsection{JARM Fingerprinting}

JARM~\cite{jarm2020} is a widely-used active TLS fingerprinting technique. It sends 10 specially
crafted \texttt{ClientHello} packets to a target server, varying parameters such as cipher suite
ordering and extension fields. Each server's TLS stack responds according to its implementation;
the pattern of selected ciphers and extensions across the 10 probes is hashed into a 62-character
fingerprint. Servers running the same TLS library and configuration produce the same JARM hash,
enabling large-scale infrastructure clustering.

\subsection{TLScope}

TLScope~\cite{theofanous2024tlscope} extends JARM-style fingerprinting with a full machine
learning pipeline. Rather than collapsing responses to a hash, TLScope treats the full per-probe
response vector as a high-dimensional feature input, preserving inter-probe correlations.
Additional features include certificate metadata, protocol capability flags, and server
behaviour under adversarial \texttt{ClientHello} configurations. The resulting classifier
achieves 91\% precision and 95\% recall for malicious server detection, and 99\% precision and
recall for botnet family identification.

\section{Threat Intelligence Feeds}

Threat intelligence feeds are curated lists of indicators of compromise (IOCs). IP-based feeds
typically assign a reputation score or binary malicious/benign label to individual IP addresses.
Feed types include:

\begin{description}
  \item[Community-reported feeds] Aggregated reports from network operators, security researchers,
    and automated systems (e.g., AbuseIPDB).
  \item[Multi-engine aggregates] Aggregation of multiple antivirus and security engines
    (e.g., VirusTotal).
  \item[Specialised IOC databases] Curated databases of malware-related indicators
    (e.g., ThreatFox).
  \item[Active scanning feeds] Feeds built from continuous internet-wide scanning
    (e.g., Censys).
  \item[Blacklists] Lightweight daily-updated IP block lists
    (e.g., Cinsscore).
  \item[Phishing URL feeds] Domain-based phishing indicators resolved to IPs
    (e.g., OpenPhish).
\end{description}

\section{Indicator Lifecycle}

The lifecycle of a threat indicator in public feeds can be characterised by three phases:
\begin{enumerate}
  \item \textbf{Pre-detection:} The threat exists but has not yet been flagged.
  \item \textbf{Active detection:} The indicator is present and the IP is flagged.
  \item \textbf{Expiry:} The indicator is removed, either because the threat has ended or
    the provider's retention policy has expired the entry.
\end{enumerate}

This thesis measures detection lag (transition from phase 1 to phase 2) and indicator lifetime
(duration of phase 2).

% ============================================================
\chapter{Related Work}
\label{ch:related_work}
% ============================================================

\section{TLScope: the Parent Project}

TLScope \cite{theofanous2024tlscope} is not a related work surveyed alongside the papers below ---
it is the \textbf{parent project} this thesis is built directly on top of, and its output is this
thesis's ground truth rather than a point of comparison. TLScope actively initiates TLS connections
to servers and sends exhaustively crafted \texttt{ClientHello} messages that progressively force
each server to select alternative cipher suites, extending JARM-style fingerprinting
(\S\ref{ch:background}) with a full ML pipeline. The system achieves 91\% precision and 95\% recall
for malicious server detection and 99\%/99\% for botnet family identification, and its output ---
groups of 1\,596--2\,728 IP addresses per batch classified as malicious --- is what
\S\ref{ch:methodology} tracks across the six public feeds.

The two projects answer different halves of the same underlying question, as
Table~\ref{tab:tlscope_netscan} makes explicit:

\begin{table}[h]
  \centering
  \caption{Division of questions between TLScope and this thesis.}
  \label{tab:tlscope_netscan}
  \begin{tabularx}{\linewidth}{Xl}
    \toprule
    \textbf{Question} & \textbf{Answered by} \\
    \midrule
    Is this server's TLS behaviour indicative of malicious hosting? & TLScope \\
    What botnet family does this server belong to?                 & TLScope \\
    How quickly do public feeds detect what TLScope already found?  & This thesis \\
    How long do providers keep an indicator flagged once detected?  & This thesis \\
    Do providers agree on the same IPs, or cover distinct subsets?  & This thesis \\
    \bottomrule
  \end{tabularx}
\end{table}

Every paper surveyed in the remainder of this chapter is instead assessed by its relation
\emph{to} TLScope, not treated as a peer of it (a curated version of this survey, cross-referenced
against this thesis's own Future Work chapter, is maintained alongside the source PDFs in
\texttt{wiki/papers/notes.md}).

\section{Malicious Domain and IP Reputation Systems}

\subsection{NOTOS}

Antonakakis et al.\ \cite{antonakakis2010notos} introduced NOTOS, a dynamic reputation system for
DNS that scores domains before they appear on static blacklists. NOTOS operates on passive DNS data
from ISP recursive resolvers, computing per-domain scores from network properties (IP diversity,
TTL volatility) and zone properties (registration age, registrar identity). It achieves a 96.8\%
true-positive rate at a 0.38\% false-positive rate and detects malicious domains weeks to months
before blacklists.

NOTOS is a foundational reference for understanding the \emph{reputation-system design space} that
public feeds implement operationally. The thesis's measured detection lag is the empirical
counterpart to NOTOS's ``lead time before blacklists''.

\subsection{EXPOSURE}

Bilge et al.\ \cite{bilge2011exposure} extended passive DNS analysis to a richer 15-feature set
covering time-based, DNS-answer-based, TTL-based, and domain-name-based features, achieving
$\sim$98\% accuracy with a C4.5 decision tree on 100 billion DNS requests. Their finding of a
1.2-day average malicious domain lifespan is directly relevant to the thesis: if TLScope-identified
IPs host similarly short-lived campaigns, public feeds may structurally miss them before evidence
accumulates. This is consistent with the thesis's own finding that 81.2\% of TLScope-identified IPs
were never flagged by any provider across the full observation window
(Chapter~\ref{ch:results}): a substantial share may simply never generate enough evidence within
a provider's detection window, regardless of how long that window is.

\subsection{MANTIS}

Deniz et al.\ \cite{deniz2025mantis} presented MANTIS, a content-agnostic system for zero-day
malicious domain detection using infrastructure reuse patterns. MANTIS builds a heterogeneous graph
over domains, IPs, ASNs, nameservers, and registrars, and applies a graph neural network to
identify domains on low-reputation hosting infrastructure. It achieves 99.7\% precision, 86.9\%
recall, and detects approximately 19\,000 new malicious domains per day --- 5$\times$ more than
VirusTotal alone.

MANTIS and this thesis address complementary lifecycle phases: MANTIS handles \emph{entry} into
threat feeds, while this thesis handles \emph{persistence and exit}. MANTIS's finding that it
detects 5$\times$ more malicious domains than VirusTotal anticipates the thesis's own measured
coverage gap: 81.2\% of TLScope-identified IPs were never flagged by any of the six tracked
providers over a one-to-two-month window (Chapter~\ref{ch:results}), a gap at least as large as
MANTIS's, though the two are measured on different units (domains vs.\ IPs) and time bases (daily
rate vs.\ full-window absence).

\section{Temporal and Longitudinal Studies of Threat Indicators}
\label{sec:rw_temporal}

Research on the temporal dimension of threat intelligence remains sparse relative to work on
detection accuracy, and where temporal results do appear they are typically a by-product of
evaluating a detector rather than the object of study. Three of the systems surveyed above report
such a by-product, and together they bound the question this thesis asks directly.

EXPOSURE \cite{bilge2011exposure} measured an average malicious-domain lifespan of 1.2 days in live
ISP deployment --- the only absolute indicator-lifetime figure in the surveyed literature, and an
order of magnitude shorter than the 8.1--32.0-day provider-side lifetimes measured here
(\S\ref{ch:results}). The two are not in conflict: EXPOSURE measures how long the malicious
\emph{infrastructure} stays active, while this thesis measures how long a \emph{provider keeps its
flag raised}, and the gap between them is precisely the interval in which a feed is reporting on
infrastructure that may already be gone. NOTOS \cite{antonakakis2010notos} and MANTIS
\cite{deniz2025mantis} both report lead time rather than lifetime --- weeks to months and days to
weeks ahead of public blacklists respectively --- but each measures that lead against blacklists as
an undifferentiated aggregate, without asking which provider supplies it or how long it persists
afterwards.

What none of these provide is a longitudinal, per-provider measurement of the full indicator
lifecycle --- entry, persistence, and expiry --- against a fixed ground-truth population tracked
over time. That is the gap this thesis addresses: rather than reporting lead time as a single
summary statistic of one detector's advantage, it tracks a fixed set of 18\,428 TLScope-identified
IPs across six providers for one to two months each and reports where each provider sits in that
lifecycle (Chapters~\ref{ch:results} and~\ref{ch:discussion}).

\section{Machine Learning for Network Security}

\subsection{Robustness Under Concept Drift}

Li et al.\ \cite{li2025conceptdrift} addressed concept drift in Windows malware detection,
using graph neural networks on control flow graphs and adversarial domain adaptation to learn
drift-invariant features. The thesis directly observes concept drift in action through temporal
tracking: as providers update their models and threat feeds, the classification of the same IP can
change over time, independently of the underlying threat state.

\subsection{Low-Quality Training Data}

Qing et al.\ \cite{qing2024rapier} presented RAPIER, a framework for robust encrypted malicious
traffic detection under low-quality training data with up to 45\% label noise. RAPIER exploits the
distributional observation that malicious traffic is scattered across feature space while benign
traffic is concentrated. This is directly applicable to the thesis: if a TLScope-identified IP is
flagged by most providers, it likely sits in a dense malicious region (high confidence); if only
one provider flags it, it likely sits in a sparse region (potential noise). The thesis's own
inter-provider agreement results confirm that the sparse case dominates in practice: 83.9\% of
ever-detected IPs are flagged by exactly one engine (Chapter~\ref{ch:results}), meaning most
provider-confirmed detections in this dataset would fall on the low-confidence, noise-prone side
of RAPIER's density distinction, and only a small minority (16.1\%) carry the multi-source
corroboration RAPIER associates with dense, high-confidence regions.

\subsection{Density-Based Robustness}

Tian et al.\ \cite{tian2025density} identified feature sparsity as a root cause of multiple
failure modes in AI-driven malware detection, including vulnerability to poisoning attacks, evasion
attacks, and concept drift. Provider coverage gaps in the thesis correspond to the sparsity
problem: TLScope-identified IPs never flagged by any public feed likely fall in sparse regions
of the threat indicator space where provider training data is thin. The scale of this coverage gap
--- 81.2\% of TLScope-identified IPs never flagged by any tracked provider (Chapter
\ref{ch:results}) --- suggests that, in this dataset, the sparse region is not a minority edge
case but the dominant regime.

\section{Internet-Wide Scanning}

\subsection{IPv6 Scanning with 6Sense}

Williams et al.\ \cite{williams2024sixsense} presented 6Sense, which uses reinforcement learning
to make IPv6 internet scanning tractable. 6Sense is a sibling project to TLScope: both perform
internet-wide active scanning, but 6Sense targets IPv6 while TLScope focuses on IPv4. The thesis
currently tracks only IPv4; 6Sense establishes the groundwork for a future IPv6 extension where
public feed coverage may lag even further.

\subsection{BGP Routing Anomaly Detection (BEAM)}

Chen et al.\ \cite{chen2024beam} presented BEAM, a semantics-aware BGP routing anomaly detection
system using AS-level role embeddings. BEAM's ``routing role consistency'' concept maps to the
thesis: each IP maintains a ``threat role'' across providers over time. Normal evolution means the
IP remains consistently flagged by the same providers; anomalies are unexpected changes in which
providers flag an IP.

\section{DNS Infrastructure and Security}

\subsection{Stale DNS Glue Records}

Zhang et al.\ \cite{zhang2024staleglue} demonstrated that 23.18\% of DNS glue records are
outdated, with over 6 million domains at risk from stale-glue domain hijacking. This directly
affects the thesis's OpenPhish data quality: the thesis converts OpenPhish URLs (domain-based) to
IP addresses via DNS resolution, and stale glue can cause incorrect IP attribution, distorting
detection lag measurements.

\subsection{DNS Resolver Vulnerabilities (ResolverFuzz)}

Zhang et al.\ \cite{zhang2024resolverfuzz} discovered 23 vulnerabilities in mainstream DNS
resolver implementations via query-response fuzzing. DNS resolver vulnerabilities in provider
infrastructure can cause providers to resolve domain-based threats to wrong IPs, appearing as
false positives or missed detections in provider comparison independently of the actual threat.

\section{Proxy-Based Abuse Detection}

Ramesh et al.\ \cite{ramesh2024calcuLatency} presented CalcuLatency, which detects proxy-enabled
abuse using cross-layer latency measurements. Different malicious infrastructure types exhibit
characteristic latency patterns. In the context of this thesis, latency-based infrastructure
clustering can explain correlated detection timing: if multiple TLScope-identified IPs share
identical latency fingerprints, they likely originate from the same hosting provider and should
exhibit similar detection lag across providers.

% ============================================================
\chapter{Methodology}
\label{ch:methodology}
% ============================================================

\section{Overview}

NetScan tracks malicious IP addresses across six public threat intelligence providers over a
one-to-two-month window. The ground truth comes from TLScope~\cite{theofanous2024tlscope}, a
machine-learning system that identifies malicious servers via active TLS fingerprinting.

The study produces two primary measurements per provider:
\begin{itemize}
  \item \textbf{Detection lag} --- how many days pass before a provider first flags a
    TLScope-identified IP.
  \item \textbf{Indicator lifetime} --- how long a provider keeps the IP classified as malicious
    once it detects it.
\end{itemize}

\section{Data Collection}

\subsection{Input Groups}

Each group consists of roughly 1\,600--2\,700 IP addresses classified as malicious by TLScope
(eight completed groups, 18\,428 IPs in total). Groups are stored in \texttt{group\_wave\_N/}
directories and each contains a CSV with columns \texttt{IP}, \texttt{Port}, and \texttt{label}.
Only rows where \texttt{label} $\neq$ \texttt{benign} are processed.

\textbf{Alignment of tracking start with TLScope classification.} Each wave is split into two or
three sub-groups, each seeded by one dated TLScope output batch, and tracking of a sub-group begins
on the date that batch was produced --- the batch directory name carries that date (e.g.\
\texttt{tlscope\_keramidas\_2026-08-05} is tracked from 2026-08-05). This alignment holds exactly
for every sub-group in waves 2 and 4--8, and within one day for waves 1 and 3, where tracking
started the day after the batch date. It is what makes RQ1 answerable at all: because tracking
starts when TLScope produces its verdict, ``cycles elapsed since the group's start date''
(\S\ref{sec:notation}) is a direct proxy for ``time elapsed since TLScope identified the server'',
which is the quantity RQ1 actually asks about. Were the two decoupled, the reported lag would
measure only when this study happened to begin looking, not how long the ecosystem took to react.

\subsection{Collection Schedule}
\label{sec:collection_schedule}

Data collection nominally runs every 3 days: for each run, \texttt{runfetch.py} queries every
configured provider for every IP in the group and writes a timestamped CSV to the wave folder.
This cadence is not perfectly uniform across the study, however. A per-group \texttt{config.json}
override (Appendix~\ref{app:cli}) allows an individual wave to run on a different cycle length, and
two of the eight waves do: wave 1 was collected on a 1-day cadence and wave 5 on a 2-day cadence,
while the remaining six waves (2--4, 6--8) use the nominal 3-day cadence. All lifetime figures
(\S\ref{sec:notation}) are computed in real days and are unaffected by this, since the cycle length
$\Delta$ is substituted per group; detection-lag figures expressed as a raw cycle count, however,
carry a different time resolution for these two waves than for the rest (\S\ref{sec:threats_to_validity}).

\subsection{Provider Availability}
\label{sec:provider_availability}

The six providers listed in \S\ref{ch:introduction} were tracked for the full study with one
exception: Censys querying (\texttt{cs\_get\_url\_report}/\texttt{process\_cs\_report} in
\texttt{vendors.py}) was disabled in the data collection tool on 2026-04-19, after wave 4 finished
and shortly before wave 5 began, and was never re-enabled. Consequently Censys results in this
thesis cover only waves 1--4 (9\,855 of the 18\,428 tracked IPs); waves 5--8 (8\,573 IPs) were
monitored by the remaining five providers only. This was a deliberate change to the collection
tool's active provider list rather than an API outage, but no rationale for it is recorded in the
codebase, and its effect on the results is addressed in \S\ref{sec:threats_to_validity}.

\subsection{Provider Query Strategies}

Table~\ref{tab:provider_queries} summarises the query type and classification basis for each
provider.

\begin{table}[h]
  \centering
  \caption{Provider query strategies and classification basis.}
  \label{tab:provider_queries}
  \begin{tabularx}{\linewidth}{lXX}
    \toprule
    \textbf{Provider} & \textbf{Query type} & \textbf{Classification basis} \\
    \midrule
    VirusTotal  & POST rescan, then GET analysis & Malicious + suspicious vote counts \\
    AbuseIPDB   & GET check (30-day window)       & \texttt{abuseConfidenceScore} threshold \\
    Censys      & GET host record                 & Presence of \texttt{c2} label \\
    ThreatFox   & POST IOC search (exact match)   & \texttt{query\_status} $\neq$ \texttt{no\_result} \\
    Cinsscore   & Local lookup in downloaded list & IP presence in CINS Army List \\
    OpenPhish   & DNS resolve + list lookup       & IP presence in resolved feed \\
    \bottomrule
  \end{tabularx}
\end{table}

\subsection{Rate Limiting}
\label{sec:rate_limiting}

VirusTotal enforces hourly and daily API quotas. A background \texttt{quota\_worker} task monitors
remaining quota every $\sim$17 minutes and cancels all fetch tasks when the daily limit is reached.
Between individual IP queries a configurable wait time of 32 seconds (\texttt{WAITING\_TIME}) is
applied to stay within rate limits. Cinsscore and OpenPhish are list-based: the full list is
downloaded once at startup and all IPs are checked locally.

\section{Analysis}

\subsection{Notation}
\label{sec:notation}

For a tracked (IP, provider) pair $(i,p)$, let $d_1 < d_2 < \dots < d_n$ be the collection dates for
$i$'s group ($\Delta = d_{t+1}-d_t = 3$ days for six of the eight waves; 1 day for wave 1 and 2 days
for wave 5, \S\ref{sec:collection_schedule}) and let
$x_t(i,p) \in \{0,1\}$ indicate whether provider $p$ classified $i$ as \emph{malicious} or
\emph{suspicious} on cycle $t$.

\textbf{Detection lag.} Let $\tau(i,p) = \min\{t : x_t(i,p) = 1\}$ be the first cycle on which $p$
detects $i$. Detection lag is measured in cycles relative to the group's start,
$\mathrm{Lag}(i,p) = \tau(i,p) - 1$; if $i$ is never detected by $p$, lag is undefined (right-censored
by the observation window rather than equal to some finite value).

\textbf{Indicator lifetime.} An indicator is considered \textbf{expired} after 2 consecutive
non-detections (\texttt{EXPIRED\_AFTER} = 2). Once first detected at $\tau(i,p)$, every subsequent
cycle --- whether a detection or a single intervening miss --- extends the indicator's lifetime by
$\Delta$; the indicator expires, and stops accumulating lifetime, on the second consecutive miss.
Concretely, if $t^*(i,p)$ is the first cycle at which two consecutive misses have occurred since
$\tau(i,p)$, the lifetime is
\[
  L(i,p) = \bigl(t^*(i,p) - \tau(i,p) - 1\bigr) \cdot \Delta ,
\]
i.e.\ the cumulative time from first detection up to, but excluding, the confirming second miss ---
a single isolated miss does not itself end the indicator, but two in a row do.

\textbf{Inter-provider agreement.} For an engine $e$ (a VirusTotal sub-engine or another provider),
let $D_e = \{(i,p) : \exists\, t,\ x_t(i,e)=1\}$ be the set of IPs it ever flags. Pairwise agreement
between engines $a$ and $b$ is the Jaccard index
\[
  J(a,b) = \frac{|D_a \cap D_b|}{|D_a \cup D_b|} .
\]

\subsection{Aggregate Statistics}

For each wave group, the following aggregate series are computed over time:
\begin{description}
  \item[\texttt{aggregate}] Total malicious votes across all IPs per time step.
  \item[\texttt{percentage}] Fraction of IPs with at least one malicious vote per time step.
  \item[\texttt{diff}] Net change in aggregate votes between consecutive time steps.
\end{description}

\subsection{Per-Provider Statistics}

For the three providers with the highest peak detection percentage, the following are generated:
\begin{itemize}
  \item Malicious percentage over time.
  \item Malicious report counts over time.
  \item Average indicator lifetime (days).
\end{itemize}

\section{Ethical Considerations}

Every provider in Chapter~\ref{ch:providers} is queried through its own documented public API or
download endpoint, within the rate limits each provider enforces (\S\ref{sec:rate_limiting}); no
provider's access controls are bypassed, and no authentication beyond a standard per-provider API
key (Appendix~\ref{app:config}) is used. NetScan does not itself connect to, probe, or exchange
traffic with the tracked malicious IPs: five of the six providers (AbuseIPDB, Censys, ThreatFox,
Cinsscore, OpenPhish) return already-collected reputation data about the target IP without NetScan
contacting that IP at all. The one exception is VirusTotal, whose rescan endpoint causes
VirusTotal's own infrastructure --- not NetScan --- to actively probe the target
(\S\ref{ch:implementation}); NetScan's role there is limited to requesting that scan and reading
back the result. No traffic to or from the tracked infrastructure is captured or decrypted by this
study. The tracked indicators are server-side infrastructure addresses drawn from an already
published upstream classification (TLScope), not end-user or personally identifiable data, and are
retained locally for research purposes only.

% ============================================================
\chapter{Threat Intelligence Providers}
\label{ch:providers}
% ============================================================

\section{VirusTotal}

VirusTotal aggregates results from over 70 antivirus engines and website scanners. For each IP,
NetScan first triggers a rescan via \texttt{POST /api/v3/ip\_addresses/\{IP\}/analyse} and then
polls the resulting analysis until \texttt{status == completed}. The \texttt{malicious +
suspicious} vote count is used as the positive signal.

VirusTotal enforces hourly and daily quotas. A background worker monitors remaining quota and
backs off on \texttt{Too Many Requests} (64 s) and \texttt{QuotaExceededError} (128 s).

\section{AbuseIPDB}

AbuseIPDB is a community-driven database of reported IPs. A single GET request with a 30-day
window returns the \texttt{abuseConfidenceScore} (0--100). NetScan applies thresholds:
$>$25 $\to$ malicious; $>$15 $\to$ suspicious; $\leq$15 $\to$ harmless.

\section{Censys}

Censys continuously scans the public IPv4 address space and maintains per-host labels including
\texttt{c2} for known C2 infrastructure. NetScan checks for \texttt{c2} labels at both the host
level and the individual service level. Any \texttt{c2} label found at either level is classified
as malicious. Censys querying was disabled in the collection tool after wave 4
(\S\ref{sec:provider_availability}); the results reported for Censys in Chapter~\ref{ch:results}
therefore cover only waves 1--4.

\section{ThreatFox}

ThreatFox (abuse.ch) is a free community platform for sharing malware IOCs. A POST request with
\texttt{exact\_match: true} is issued; any \texttt{query\_status} other than \texttt{no\_result}
is classified as malicious.

\section{Cinsscore}

The CINS Army List is a public text file updated continuously. NetScan downloads the full list
once at startup and checks all IPs locally. This is the simplest provider integration: a binary
in-list/not-in-list classification with no per-IP API overhead.

\section{OpenPhish}

OpenPhish provides a domain-based phishing URL feed. Because the feed contains URLs rather than
IP addresses, NetScan extracts hostnames and resolves them to IPs via DNS at startup. An IP found
in the resolved set is classified as malicious.

The DNS resolution step introduces noise: domains that rotate IPs (fast-flux, CDN-backed phishing)
may produce stale associations. Stale DNS glue records \cite{zhang2024staleglue} are an additional
source of incorrect IP attribution.

% ============================================================
\chapter{Implementation}
\label{ch:implementation}
% ============================================================

\section{Repository Structure}

\begin{table}[h]
  \centering
  \caption{NetScan repository files and their purpose.}
  \label{tab:repo}
  \begin{tabularx}{\linewidth}{lX}
    \toprule
    \textbf{File} & \textbf{Purpose} \\
    \midrule
    \texttt{runfetch.py}       & Async data collection from all providers \\
    \texttt{runanalysis.py}    & Analysis pipeline and graph generation \\
    \texttt{vendors.py}        & Provider API client functions and report processors \\
    \texttt{available\_fields.py} & Utility to inspect CSV output field schema \\
    \texttt{.env\_ex}          & Example environment / API key configuration \\
    \bottomrule
  \end{tabularx}
\end{table}

\section{Fetch Pipeline}

VirusTotal's per-IP analysis is asynchronous on VirusTotal's own side: a submitted rescan must be
polled until its status transitions to \texttt{completed}, and each poll iteration itself waits a
multiple of \texttt{WAITING\_TIME} (\S\ref{ch:providers}, VirusTotal). An earlier, single-pool
design queried every provider for a given IP from the same worker, so that worker sat blocked on
VirusTotal's polling loop before it could move on to the next IP, even though the other providers'
queries for that same IP had already completed in a fraction of the time. The pipeline is
consequently split into two cooperating pools connected by a second, internal queue:

\begin{enumerate}
  \item At startup, Cinsscore and OpenPhish lists are downloaded once and held in memory.
  \item For each input CSV in the folder, a new output file is created.
  \item A pool of \texttt{-w} \textbf{producer} workers pulls IPs from the input queue. For each IP,
    a producer submits the VirusTotal rescan \emph{without waiting for it to complete}, then queries
    every other configured provider (AbuseIPDB, ThreatFox, Cinsscore, OpenPhish) synchronously. The
    partial result --- including the VirusTotal analysis id --- is handed off to a second,
    VirusTotal-result queue.
  \item A pool of \texttt{-vw} \textbf{consumer} workers pulls partial results from that second
    queue, polls the corresponding VirusTotal analysis until it reports \texttt{completed}, merges
    the result into the row, and writes the completed row to the output CSV.
  \item A separate \texttt{quota\_worker} background task monitors VirusTotal's daily quota and
    cancels every producer and consumer task if it is exhausted.
  \item Producers sleep 32 seconds between IPs to respect per-provider rate limits; the VirusTotal
    polling loop inside each consumer applies its own independent backoff.
\end{enumerate}

This decouples the dominant source of per-IP latency --- VirusTotal polling --- from retrieval of
every other provider's data: while a consumer waits on VirusTotal for one IP, producers continue
submitting rescans and fetching AbuseIPDB/ThreatFox/Cinsscore/OpenPhish data for subsequent IPs,
rather than the whole pipeline stalling on a single in-flight VirusTotal analysis. File I/O (CSV
input/output, per-wave output folder creation) is performed asynchronously via the \texttt{anyio}
library so that large file reads and writes do not block the same event loop that drives the
provider queries.

\section{Fetch Output Format}

Each report CSV contains per-IP columns including \texttt{total\_votes} (aggregate vote counts),
\texttt{report} (per-provider binary classifications), and raw provider response fields
(\texttt{vt\_response}, \texttt{ai\_response}, \texttt{tf\_response}, etc.).

\section{Analysis Output}

Running \texttt{runanalysis.py} on a wave folder writes the following files to an output
subdirectory (\texttt{group\_analysis\_out\_dir/} by default) inside that folder:

\begin{table}[h]
  \centering
  \caption{Files produced by the analysis pipeline.}
  \label{tab:analysis_output}
  \begin{tabularx}{\linewidth}{lX}
    \toprule
    \textbf{File} & \textbf{Content} \\
    \midrule
    \texttt{out\_logs.txt}       & Run arguments and all computed statistics (Total Stats, Report Stats) \\
    \texttt{all\_zero.csv}       & IPs with zero malicious votes from any provider \\
    \texttt{non\_zero.csv}       & IPs detected by at least one provider at least once \\
    \texttt{ip\_port\_series.csv}  & Per-host vote history (\texttt{IP}, \texttt{Port}, \texttt{Group}, \texttt{Date}, votes) \\
    \texttt{total\_series.csv}   & Aggregate/diff/percentage series, one row per collection cycle \\
    \texttt{engine\_series.csv}  & Per-engine count/percentage series, one row per engine per cycle \\
    \texttt{engine\_summary.csv} & Per-engine average lifetime, peak detection percentage, and
      ever-detected IP count (\S\ref{sec:agreement_vs_chance}) \\
    \texttt{group\_available\_fields.json} & Merged schema of every field across all provider
      responses in the wave's report files (per-engine field groups collapsed to one
      representative sample), for schema inspection without hand-parsing raw CSVs \\
    \texttt{figure\_*.png}       & Generated graphs (aggregate, diff, per-host, top-engine detail) \\
    \bottomrule
  \end{tabularx}
\end{table}

All numerical results in Chapter~\ref{ch:results} are drawn directly from these files.

% ============================================================
\chapter{Results}
\label{ch:results}
% ============================================================

This chapter presents empirical results from eight completed collection waves
(\texttt{group\_wave\_1} through \texttt{group\_wave\_8}), covering 18\,428 TLScope-classified IP
observations tracked on a predominantly three-day cadence (\S\ref{sec:collection_schedule}) over
one-to-two-month windows between October 2025 and
September 2026. Per-wave figures and raw statistics are available under each wave's
\texttt{group\_analysis\_out\_dir/}; this chapter reports the aggregate and per-wave numbers drawn
from them. The methodological caveats that bound how these numbers should be read --- ground-truth
imperfection, sampling granularity, and selection bias among them --- are addressed together in
\S\ref{sec:threats_to_validity} rather than repeated at each figure.

\section{Detection Rate Overview}

Table~\ref{tab:detection_overview} summarises, for each wave, how many TLScope-identified IPs were
ever flagged by at least one tracked provider versus never flagged during the entire observation
window.

\begin{table}[h]
  \centering
  \caption{Detection outcome per wave: IPs never flagged by any provider versus ever detected.}
  \label{tab:detection_overview}
  \begin{tabular}{lrrrr}
    \toprule
    \textbf{Wave} & \textbf{Total IPs} & \textbf{Never detected} & \textbf{Ever detected} & \textbf{Never detected \%} \\
    \midrule
    1 & 2\,591 & 1\,900 & 691 & 73.3\% \\
    2 & 2\,062 & 1\,507 & 555 & 73.1\% \\
    3 & 2\,728 & 2\,188 & 540 & 80.2\% \\
    4 & 2\,474 & 2\,144 & 330 & 86.7\% \\
    5 & 1\,596 & 1\,390 & 206 & 87.1\% \\
    6 & 2\,470 & 2\,079 & 391 & 84.2\% \\
    7 & 2\,100 & 1\,772 & 328 & 84.4\% \\
    8 & 2\,407 & 1\,987 & 420 & 82.6\% \\
    \midrule
    \textbf{Total} & \textbf{18\,428} & \textbf{14\,967} & \textbf{3\,461} & \textbf{81.2\%} \\
    \bottomrule
  \end{tabular}
\end{table}

Across all eight waves, \textbf{81.2\% of TLScope-identified malicious IPs (14\,967 of 18\,428)
were never flagged by any of the tracked providers} at any point during their observation window.
This is the headline result of the thesis: the large majority of servers that active TLS
fingerprinting identifies as malicious never surface in public threat intelligence feeds at all,
independent of detection lag.

The never-detected rate ranges from 73.1\% to 87.1\% across waves and is not flat: it rises with
wave order from $\sim$73\% in waves 1--2 to a 82--87\% band from wave 4 onward. The correlation
with wave order is substantial but, at eight waves, falls just short of conventional significance
(Pearson $r = 0.704$, $p = 0.051$, $n = 8$) --- it should be read as a consistent direction of
travel, not as an established effect. Part of the rise plausibly coincides with the Censys
deactivation after wave 4 (\S\ref{sec:provider_availability}): the mean never-detected rate is
78.3\% for waves 1--4 (Censys active) versus 84.5\% for waves 5--8 (Censys inactive), a difference
that likewise does not reach significance with four waves on each side ($t = -1.85$, $p = 0.114$).
Neither test settles whether the trend reflects a genuine ecosystem shift or the loss of a
provider, and the honest summary is the weaker one: across this study's window, coverage shows no
sign of improving, and the apparent decline is confounded with a change in what was being
measured. Distinguishing the two requires either more waves or the restoration of Censys
(\S\ref{ch:future_work}).

\section{Detection Lag}

Because collection runs once per cycle rather than continuously, detection lag is measured in units
of one collection cycle rather than in continuous days. A cycle is 3 days for six of the eight
waves, but only 1 day for wave 1 and 2 days for wave 5 (\S\ref{sec:collection_schedule}); ``zero
lag'' is consequently a stricter, same-day claim for wave 1 than the up-to-3-day claim it represents
for most other waves, and the figures below mix these two resolutions rather than reporting a
single uniform unit. Of the 3\,461 IPs that were eventually detected at least
once, 1\,523 (44.0\%) were flagged on the very first collection cycle (zero lag relative to the
group's start date), while 1\,410 (40.7\%) were not detected on the first cycle but were flagged on
a later one. The remaining IPs correspond to groups whose first available sample fell outside the
expected collection schedule and are excluded from this lag breakdown.

The immediate-detection fraction varies substantially by wave, from 27.6\% (wave 8) to 70.4\%
(wave 5), with no consistent trend across waves (Pearson $r = -0.03$ against wave order). This
variation is consistent with detection lag depending heavily on which specific hosting
infrastructure and campaign type a given wave's IPs belong to, rather than on a fixed system-wide
latency --- though the wave-1 figure in particular should be read against its finer 1-day cadence
rather than compared at face value to waves on a 3-day cycle.

\section{Indicator Lifetime}

Per-engine lifetime and peak-detection statistics are computed for every distinct engine that
appears in a provider's report --- for VirusTotal specifically, this means each of the
third-party antivirus and blocklist engines VirusTotal aggregates internally (e.g.\ \emph{Criminal
IP}, \emph{SOCRadar}, \emph{BitDefender}), rather than VirusTotal as a single aggregate. The
analysis reports the five highest peak-detection engines per wave (\texttt{-t 5},
Appendix~\ref{app:cli}). Exactly two engines appear in that set in every one of the eight waves ---
\emph{Criminal IP} and \emph{alphaMountain.ai} --- and Table~\ref{tab:top_engine} tracks both,
because they are consistent in opposite ways.

\begin{table}[h]
  \centering
  \caption{Average indicator lifetime and peak detection percentage for the two engines present
    among the five highest peak-detection engines in all eight waves.}
  \label{tab:top_engine}
  \begin{tabular}{lrrrr}
    \toprule
    & \multicolumn{2}{c}{\textbf{Criminal IP}} & \multicolumn{2}{c}{\textbf{alphaMountain.ai}} \\
    \cmidrule(lr){2-3}\cmidrule(lr){4-5}
    \textbf{Wave} & \textbf{Lifetime (days)} & \textbf{Max det.\ \%}
                  & \textbf{Lifetime (days)} & \textbf{Max det.\ \%} \\
    \midrule
    1 & 8.5  & 17.9\% & 7.5  & 0.5\% \\
    2 & 22.1 & 18.7\% & 33.6 & 1.7\% \\
    3 & 25.5 & 12.8\% & 24.6 & 1.7\% \\
    4 & 32.0 & 8.0\%  & 20.0 & 3.0\% \\
    5 & 11.1 & 5.0\%  & 37.0 & 3.0\% \\
    6 & 11.3 & 5.5\%  & 33.3 & 1.5\% \\
    7 & 8.4  & 4.0\%  & 36.0 & 2.3\% \\
    8 & 8.1  & 4.3\%  & 30.4 & 2.3\% \\
    \bottomrule
  \end{tabular}
\end{table}

The two behave as near-opposites, and the contrast is the substantive result here. \emph{Criminal
IP} is the \emph{broad} detector: it reaches by far the highest detection share of any engine
(17.9\% in wave 1) but holds its indicators only briefly and erratically --- 8.1 to 32.0 days,
almost a 4$\times$ spread --- and its peak share decays steadily across the study, from $\sim$18\%
in waves 1--2 to $\sim$4--8\% from wave 4 onward. \emph{alphaMountain.ai} is the \emph{persistent}
detector: it never exceeds 3.0\% of a wave, but once it flags an IP it holds that flag for 20--37
days in every wave after the first. Breadth and persistence are therefore not merely uncorrelated
in this dataset; among the only two engines present throughout, they are inversely related.

The remaining engines turn over. \emph{SOCRadar} appears in seven of eight waves (1--7; average
lifetime 7.9--32.4 days); \emph{AbuseIPDB} (waves 2, 6, 7) and \emph{Gridinsoft} (waves 6--8) in
three each; \emph{ThreatFox} (1, 3), \emph{Fortinet} (3, 4) and \emph{ESET} (5, 8) in two each; and
\emph{Trustwave}, \emph{CRDF}, \emph{MalwareURL}, \emph{Seclookup} and \emph{Chong Lua Dao} in one
wave apiece. Thirteen distinct engines occupy the 40 available slots across the study, so which
engines are actively contributing detections is itself unstable over an eleven-month window --- a
point that matters for anyone treating a specific engine's coverage as a fixed property.

\section{Provider-Level Coverage}
\label{sec:provider_coverage}

The results above operate mostly at the level of VirusTotal's individual aggregated engines, since
those are what compete for the highest per-wave detection percentages. It is also useful to ask the
question the six named providers in \S\ref{ch:introduction} were originally chosen to answer
directly: of the 3\,461 ever-detected IPs, how many did each \emph{standalone} provider (i.e.\ every
provider other than VirusTotal, whose own contribution is the union of the many engines already
covered above) account for on its own? Table~\ref{tab:provider_coverage} answers this by summing
each provider's ever-detected count (\texttt{EverDetectedCount} in \texttt{engine\_summary.csv})
across every wave in which it was active.

\begin{table}[h]
  \centering
  \caption{Aggregate ever-detected IP count for each standalone (non-VirusTotal) provider, summed
    across every wave in which it was active. Censys's percentage is of the 9\,855 IPs in waves
    1--4, the only waves in which it was queried (\S\ref{sec:provider_availability}); all other
    percentages are of the full 18\,428-IP population.}
  \label{tab:provider_coverage}
  \begin{tabular}{lrrr}
    \toprule
    \textbf{Provider} & \textbf{Ever-detected IPs} & \textbf{\% of applicable population} & \textbf{Avg.\ lifetime (days)} \\
    \midrule
    AbuseIPDB  & 180 & 0.98\%                  & 23.8 \\
    ThreatFox  & 259 & 1.40\%                  & 4.7  \\
    Censys     & 16  & 0.16\% (waves 1--4 only) & 3.4  \\
    OpenPhish  & 2   & 0.01\%                  & 7.5  \\
    Cinsscore  & 0   & 0.00\%                  & ---  \\
    \bottomrule
  \end{tabular}
\end{table}

Two results stand out. First, \textbf{Cinsscore never flagged a single one of the 18\,428 tracked
IPs, in any of the eight waves}: as a lightweight, list-only blacklist (\S\ref{ch:providers}), it
contributes nothing to this dataset's coverage of TLScope-identified infrastructure. Second,
\textbf{OpenPhish flagged only 2 IPs in total}, both in wave 3. This is best read as a construct
mismatch rather than a coverage failure: OpenPhish is a phishing-\emph{URL} feed resolved to IPs by
DNS, while TLScope's ground truth targets TLS-fingerprinted C2/botnet server infrastructure
(\S\ref{ch:background}) --- largely disjoint threat categories that happen to share the same
IP-address representation. By contrast, ThreatFox and AbuseIPDB are genuine, if modest, contributors
(1.40\% and 0.98\% of all tracked IPs respectively), consistent with their occasional appearances in
the top-3 tables above. Censys's 0.16\% figure cannot be compared on equal terms with the others,
since it reflects only half the study's waves and roughly half its IPs; whether Censys would have
continued at a similarly low rate, or behaved differently, for waves 5--8 is unknown
(\S\ref{sec:threats_to_validity}). Together, these five providers account for 457 of the 3\,461
ever-detected IPs (13.2\%); the remaining 86.8\% of detections are attributable to VirusTotal's
aggregated third-party engines discussed above, confirming that VirusTotal --- not the five
standalone providers --- is the dominant source of whatever coverage this dataset's six-provider
design does achieve.

\section{Inter-Provider Agreement}

Table~\ref{tab:agreement} aggregates, across all eight waves, how many detected engines flagged
each of the 3\,461 ever-detected IPs.

\begin{table}[h]
  \centering
  \caption{Number of detecting engines per ever-detected IP, aggregated across all eight waves.}
  \label{tab:agreement}
  \begin{tabular}{lrr}
    \toprule
    \textbf{Engines agreeing} & \textbf{IP count} & \textbf{\% of ever-detected} \\
    \midrule
    Exactly 1  & 2\,903 & 83.9\% \\
    Exactly 2  & 338    & 9.8\%  \\
    Exactly 3  & 99     & 2.9\%  \\
    4 or more  & 121    & 3.5\%  \\
    \midrule
    \textbf{Total} & \textbf{3\,461} & \textbf{100.0\%} \\
    \bottomrule
  \end{tabular}
\end{table}

Agreement between engines is low: 83.9\% of ever-detected IPs are flagged by exactly one engine,
and fewer than 4\% are flagged by four or more. This confirms that different engines cover largely
distinct, only lightly overlapping subsets of the malicious IP population, directly answering RQ3.

One clear exception recurs across every wave: the pairwise Jaccard similarity between
\emph{BitDefender} and \emph{G-Data} is exactly 1.0 in all eight waves --- every IP one flags, the
other flags too, with no exceptions in any wave. This is consistent with the two commercial
products sharing the same underlying detection engine, which is a documented licensing
relationship between the two vendors, rather than coincidental correlated detection.

\subsection{Agreement Versus a Chance Baseline}
\label{sec:agreement_vs_chance}

The low overlap in Table~\ref{tab:agreement} is consistent with two different explanations that a
purely descriptive count cannot distinguish: either individual engines simply detect few IPs each
(so singleton detections would dominate even if every engine's hits were an independent random
draw), or engines are additionally \emph{specialising} --- systematically detecting different IPs
from one another, beyond what their individual detection rates alone would predict. To separate
these, for every wave and every engine $e$ with at least one detection, the marginal probability
$p_e$ that $e$ ever flags a given ever-detected IP is estimated as $p_e = |D_e| / N$, where $N$ is
that wave's ever-detected count (Table~\ref{tab:detection_overview}) and $D_e$ is as defined in
\S\ref{sec:notation}. Treating each engine as an independent Bernoulli$(p_e)$ draw per IP, the
Poisson binomial distribution over the number of engines flagging a given IP gives a
chance-baseline prediction for the fraction of IPs flagged by exactly one engine, which can be
compared against the value actually observed in that wave.

\begin{table}[h]
  \centering
  \caption{Observed versus chance-predicted fraction of ever-detected IPs flagged by exactly one
    engine, per wave. The chance prediction uses each engine's own marginal detection rate for that
    wave (Poisson binomial model, \S\ref{sec:agreement_vs_chance}); observed exceeding predicted
    indicates specialisation beyond what individually low detection rates alone would produce.}
  \label{tab:agreement_chance}
  \begin{tabular}{lrrrr}
    \toprule
    \textbf{Wave} & \textbf{Engines} & \textbf{Observed} & \textbf{Chance-predicted} & \textbf{Ratio} \\
    \midrule
    1 & 21 & 86.8\% & 59.1\% & 1.47$\times$ \\
    2 & 26 & 74.8\% & 43.1\% & 1.74$\times$ \\
    3 & 36 & 81.1\% & 37.0\% & 2.19$\times$ \\
    4 & 31 & 86.7\% & 43.2\% & 2.00$\times$ \\
    5 & 23 & 91.3\% & 49.0\% & 1.86$\times$ \\
    6 & 33 & 88.7\% & 50.5\% & 1.76$\times$ \\
    7 & 33 & 80.8\% & 39.7\% & 2.03$\times$ \\
    8 & 31 & 86.7\% & 47.4\% & 1.83$\times$ \\
    \bottomrule
  \end{tabular}
\end{table}

In every one of the eight waves, the observed exactly-one-engine fraction exceeds the chance
prediction by a wide and consistent margin --- 1.47$\times$ to 2.19$\times$ --- rather than tracking
it closely, as it would if singleton detections were merely an artefact of each engine's own low
individual detection rate. This is direct evidence that engines are not drawing independently from
a shared population of malicious IPs: even after accounting for how rarely each engine detects
\emph{anything}, IPs that one engine detects are still less likely than chance to also be detected
by another. Combined with the near-total BitDefender/G-Data overlap immediately above, the two
findings together indicate that inter-provider disagreement in this dataset is not simple
sparsity, but reflects real specialisation across largely non-redundant detection surfaces, with
shared-engine licensing relationships (as with BitDefender/G-Data) the one identifiable exception.
The independence assumption underlying the chance baseline is itself a simplification --- it treats
every ever-detected IP as facing identical per-engine marginal probabilities, ignoring any
correlation induced by shared hosting infrastructure across IPs within a wave --- so the reported
ratios should be read as a lower bound on the true degree of specialisation, not a precise estimate.

\section{Temporal Evolution}

Wave 4's aggregate series (14 collection points over 39 days) is representative of the general
pattern observed across waves: the percentage of IPs with at least one malicious vote stays within
a narrow band (9.6\%--12.4\%) for the full duration, with no sustained upward or downward trend
(Figure~\ref{fig:wave4_percentage}). The step-to-step net change (\texttt{diff}) oscillates between
roughly $-14$ and $+22$ aggregate votes per cycle, with positive and negative contributions of
comparable magnitude at every step (Figure~\ref{fig:wave4_diff}) --- new detections and lost
detections roughly balance rather than accumulating in one direction. This suggests that, over a
one-to-two-month window, provider coverage of a given wave's IPs reaches a rough equilibrium
quickly rather than continuing to climb as feeds ``catch up'' over time.

\begin{figure}[h]
  \centering
  \includegraphics[width=0.85\linewidth]{wave4_percentage.png}
  \caption{Percentage of wave 4 IPs with at least one malicious vote, per collection cycle.}
  \label{fig:wave4_percentage}
\end{figure}

\begin{figure}[h]
  \centering
  \includegraphics[width=0.85\linewidth]{wave4_diff.png}
  \caption{Net, positive, and negative components of the step-to-step aggregate vote change for
    wave 4.}
  \label{fig:wave4_diff}
\end{figure}

% ============================================================
\chapter{Discussion}
\label{ch:discussion}
% ============================================================

\section{Answers to the Research Questions}
\label{sec:rq_answers}

Before interpreting the results further, this section states directly what the measurements in
Chapter~\ref{ch:results} say about each of the three research questions posed in
\S\ref{ch:introduction}. A precondition applies to all three: they are answerable only for the
18.8\% of tracked IPs a provider ever flagged, since the remaining 81.2\% produce no detection
event to measure lag, lifetime, or agreement on. That censoring is itself the single largest
finding, and it bounds every answer below.

\textbf{RQ1 (detection lag) --- bimodal, and mostly immediate when it happens at all.} Of the
3\,461 ever-detected IPs, 44.0\% were flagged on the first collection cycle and 40.7\% only on a
later one (\S\ref{ch:results}); the remainder fall outside the schedule and are excluded. Providers
therefore do not converge on a characteristic reaction time: an IP is either already known when
TLScope surfaces it, or it waits. The immediate-detection share swings from 27.6\% to 70.4\% across
waves with no trend over time ($r = -0.01$), which points to the composition of each wave's
infrastructure rather than to any fixed ecosystem latency. The lag is measured in collection cycles
rather than days, and those cycles are not the same length in every wave
(\S\ref{sec:collection_schedule}), so this answer is directional rather than a precise interval.

\textbf{RQ2 (indicator persistence) --- highly variable, and inversely related to how widely an
engine detects.} The two engines present in every wave's top five answer this most directly
(Table~\ref{tab:top_engine}): \emph{Criminal IP} detects broadly (peaking at 17.9\% of a wave) but
holds indicators only 8.1--32.0 days, an almost 4$\times$ spread, with its share decaying across the
study; \emph{alphaMountain.ai} never exceeds 3.0\% of a wave yet retains flags for 20--37 days in
every wave after the first. Across the five standalone providers, average lifetimes span 3.4 days
(Censys) to 23.8 days (AbuseIPDB) (Table~\ref{tab:provider_coverage}). No engine is simultaneously
broad and persistent; breadth and persistence trade off against each other, and the absolute
lifetime also tracks the wave --- consistent with retention being driven both by an engine's own
policy and by what a given batch of infrastructure was doing.

\textbf{RQ3 (inter-provider agreement) --- near-absent, and by specialisation rather than by
scarcity.} 83.9\% of ever-detected IPs are flagged by exactly one engine and under 4\% by four or
more (Table~\ref{tab:agreement}). The chance-baseline comparison
(\S\ref{sec:agreement_vs_chance}) shows this is not merely an artefact of each engine detecting
little: the observed single-engine share exceeds what independent draws at each engine's own
detection rate would predict, by 1.47$\times$--2.19$\times$ in every wave. Engines are covering
systematically different subsets. The one exception is BitDefender and G-Data, identical in all
eight waves (Jaccard = 1.0), which reflects a shared underlying engine rather than agreement
between independent assessments.

\section{Coverage Gap between TLScope and Public Feeds}

The 81.2\% never-detected rate (Table~\ref{tab:detection_overview}) quantifies the gap between
what TLScope identifies through active TLS scanning and what public feeds cover: four out of five
TLScope-flagged malicious servers never appear in any of the six tracked providers, at any
granularity, over a one-to-two-month window. This gap has been hypothesised by prior work (e.g.,
MANTIS \cite{deniz2025mantis} detects 5$\times$ more malicious domains than VirusTotal alone), and
the measurement here is consistent in direction and, if anything, larger in magnitude --- though
the two are not directly comparable, since MANTIS measures domains and daily detection rate,
while this thesis measures IPs never detected across an entire multi-week window.

\section{Provider-Specific Behaviours}

The per-engine results (Table~\ref{tab:top_engine}) show that no single engine dominates
consistently across both detection lag and lifetime. Only two engines --- \emph{Criminal IP} and
\emph{alphaMountain.ai} --- appear in every wave's top five, and they split the two properties
between them rather than combining them: the broad detector is the impersistent one and the
persistent detector is the narrow one. \emph{Criminal IP}'s average lifetime also varies by nearly
4$\times$ across waves (8.1--32.0 days) while its peak detection percentage trends downward in
later waves --- evidence against a single fixed detection-lag or lifetime profile per engine, and
more consistent with behaviour that depends on the specific hosting infrastructure and campaign
type present in a given wave. The absence of Censys, Cinsscore, and OpenPhish from the top-engine
tables across all eight waves has two distinct causes rather than one, as Table
\ref{tab:provider_coverage} (\S\ref{sec:provider_coverage}) shows: Cinsscore and OpenPhish were
tracked in every run but genuinely detected almost nothing (0 and 2 IPs respectively, out of
18\,428), whereas Censys was tracked in only half the waves before being deactivated in the
collection tool (\S\ref{sec:provider_availability}) and so cannot be judged against the other four
on equal terms. In all three cases, though, their practical contribution to catching
TLScope-identified IPs in this dataset is small relative to VirusTotal's aggregated third-party
engines, which account for 86.8\% of every detection made (\S\ref{sec:provider_coverage}). The
chance-baseline comparison in
\S\ref{sec:agreement_vs_chance} reinforces this picture quantitatively: the observed single-engine
detection rate exceeds what each engine's own marginal detection rate alone would predict by
1.47$\times$--2.19$\times$ in every wave, so the low agreement in Table~\ref{tab:agreement} reflects
engines specialising in largely distinct subsets of the malicious IP population, not merely each
engine independently detecting few IPs.

\section{Indicator Expiry Mechanisms}

Indicators expire for fundamentally different reasons across providers:
\begin{enumerate}
  \item The underlying threat has genuinely ended (infrastructure taken down).
  \item Provider retention policy removed the entry.
  \item Classification model updated (concept drift \cite{li2025conceptdrift}).
\end{enumerate}
The thesis's temporal tracking distinguishes the first two cases empirically but cannot directly
identify the third.

\section{Threats to Validity}
\label{sec:threats_to_validity}

\subsection{Internal Validity}

\textbf{Ground-truth imperfection.} TLScope reports 91\% precision for malicious server detection
(\S\ref{ch:background}), so up to $\sim$9\% of the 18\,428 tracked IPs could be TLScope false
positives --- benign servers that no provider should ever flag. Any such IP correctly contributes to
the 81.2\% never-detected figure for reasons unrelated to feed coverage, meaning that figure is an
upper bound on the true ecosystem coverage gap, not an unbiased point estimate. Because a benign IP
could still be flagged by a provider for an unrelated reason (e.g.\ shared hosting with a genuinely
malicious neighbour), the direction of the resulting bias on the never-detected rate specifically is
downward (toward overstating the gap), while its effect on detection lag and lifetime is unsigned.

\textbf{Sampling granularity.} The collection cadence caps the resolution of every lag and lifetime
measurement, and this cadence is not uniform across the study: $\Delta = 3$ days for six of the
eight waves, but $\Delta = 1$ day for wave 1 and $\Delta = 2$ days for wave 5
(\S\ref{sec:collection_schedule}). The 44.0\% ``zero lag'' figure (\S\ref{ch:results}) is therefore
an upper bound on true same-day detection for the six 3-day waves --- some of those IPs may
actually have been detected 1--2 days after TLScope but before the next scheduled sample, and would
be indistinguishable from a same-cycle detection --- while for wave 1 specifically, ``zero lag''
is a same-day claim with no such slack. Aggregating cycle counts across waves with different
$\Delta$ therefore mixes two different time resolutions into one headline percentage, which should
be read as directionally informative rather than as a single precisely-defined quantity.

\textbf{Inconsistent provider coverage.} Censys querying was deactivated in the collection tool
after wave 4 (\S\ref{sec:provider_availability}) and never restored, so the ``six providers''
framing used throughout this thesis holds fully only for waves 1--4 (9\,855 IPs); waves 5--8
(8\,573 IPs) were monitored by five providers. Because Censys contributed only 16 detections in the
waves it was active for (\S\ref{sec:provider_coverage}), its removal is unlikely to explain most of
the coverage gap on its own, but it mechanically raises the never-detected rate for every wave after
it by removing one chance at detection per IP, and is a plausible partial contributor to the
increasing never-detected trend noted in \S\ref{ch:results}. No equivalent gap exists for the other
five providers, which returned a live (possibly zero-valued) result for every IP in every wave they
were configured for.

\subsection{Construct Validity}

\textbf{IP churn.} An IP address is not a stable proxy for a single threat: hosting providers
reassign addresses, and a formerly malicious IP can turn benign (or vice versa) within the
one-to-two-month observation window without any change of classification target. Both TLScope's
ground-truth snapshot and every provider's periodic classification could be tracking different
physical infrastructure occupying the same address at different points in time, which this thesis's
IP-level tracking cannot distinguish from a genuine detection or expiry event.

\textbf{No benign control set.} All 18\,428 tracked IPs are drawn from a single malicious-only
ground truth; this thesis measures each provider's recall-like coverage of that population but does
not estimate false-positive rate or precision, since no known-benign IPs are tracked in parallel.

\textbf{Provider-internal drift.} The set of individual antivirus/blocklist engines VirusTotal
aggregates is decided by VirusTotal, not by this thesis, and can change during the collection
window --- an engine can be added, removed, or re-licensed, as the identical BitDefender/G-Data
detections in \S\ref{ch:results} illustrate. Such a change would appear in the data as a shift in
per-engine coverage with no corresponding shift in the underlying threat.

\subsection{External Validity}

\textbf{Detection-methodology selection bias.} TLScope selects on TLS handshake behaviour. Malicious
infrastructure that TLScope's fingerprinting approach does not surface (e.g.\ non-TLS abuse, or
servers that closely mimic benign TLS stacks) is structurally invisible to this study. The measured
coverage gap therefore characterises ecosystem response specifically to the subset of malicious
infrastructure that active TLS fingerprinting happens to find --- which may be exactly the subset
that signature- or report-based feeds (e.g.\ ThreatFox, AbuseIPDB) are least tuned to catch. This
concern is compounded by the study being IPv4-only, since IPv6 infrastructure is entirely out of
scope (\S\ref{ch:future_work}).

\textbf{Single vantage point.} All provider queries are issued from one network location using
standard REST APIs and a single API key per provider. Provider responses could differ by client
identity, subscription tier, or geographic origin of the query, none of which this thesis observes
or controls for.

% ============================================================
\chapter{Future Work}
\label{ch:future_work}
% ============================================================

\begin{description}
  \item[IPv6 extension] The study covers IPv4 only. As IPv6 adoption grows and malicious
    infrastructure migrates to IPv6, extending tracking to IPv6 addresses using 6Sense-style
    scanning \cite{williams2024sixsense} would be valuable.
  \item[Restore or replace Censys] Censys was deactivated in the collection tool after wave 4
    (\S\ref{sec:provider_availability}) without a recorded rationale. Re-enabling it, or replacing
    it with a comparable active-scanning feed, would restore genuine six-provider coverage for
    future waves and remove the confound identified in \S\ref{sec:threats_to_validity}.
  \item[Provider health monitoring] Neither the Censys deactivation nor a transient per-IP failure
    (\texttt{cs\_get\_url\_report}/\texttt{process\_cs\_report} in \texttt{vendors.py} silently
    drop the field on any exception, with no fallback value) would be visible without inspecting
    the output schema by hand, as this thesis did. Logging or alerting when a configured provider
    stops returning results for an extended period would catch this class of issue as it happens
    rather than after the fact.
  \item[Harmonised sampling cadence] Waves 1 and 5 ran on 1-day and 2-day cadences respectively
    while the rest used 3 days (\S\ref{sec:collection_schedule}). Standardising on a single cadence
    for all future waves would make cross-wave detection-lag comparisons directly commensurable
    without the caveats in \S\ref{sec:threats_to_validity}.
  \item[Additional providers] Adding more providers (e.g., Shodan, IPinfo, AlienVault OTX) would
    give a more complete view of the threat intelligence ecosystem.
  \item[Automated provider profiling] Using MANTIS-style infrastructure graph features
    \cite{deniz2025mantis} to predict which providers are likely to detect a given IP, and when.
  \item[AS-level feature fusion] \S\ref{ch:related_work} draws a direct analogy between BEAM's
    \cite{chen2024beam} per-AS ``routing role'' embeddings and the notion of a per-IP ``threat
    role'' across providers, but this thesis never tests it: adding AS/BGP-derived features
    (announcing AS, role stability, peer diversity) to the per-engine analysis in
    \S\ref{ch:results} could test whether the specialisation observed in
    \S\ref{sec:agreement_vs_chance} correlates with hosting-infrastructure role rather than being
    engine-idiosyncratic.
  \item[Latency-based infrastructure clustering] Similarly, CalcuLatency's \cite{ramesh2024calcuLatency}
    cross-layer latency fingerprinting is discussed in \S\ref{ch:related_work} as a way to detect
    IPs sharing the same hosting provider, but this thesis does not apply it. Clustering
    TLScope-identified IPs by latency fingerprint would let a future wave test directly whether
    IPs in the same cluster share detection lag and lifetime, rather than only observing that
    lifetime varies by wave (\S\ref{ch:results}) without knowing why.
  \item[Density-aware validation of the coverage gap] \S\ref{ch:related_work} interprets the
    81.2\% never-detected rate as consistent with TLScope-identified IPs occupying sparse regions
    of provider feature space, drawing on Tian et al.'s \cite{tian2025density} finding that sparse
    regions are where malware classifiers fail hardest. That interpretation is not tested in this
    thesis: applying their density-compression measurement to whatever feature representation a
    given provider's classifier uses (where accessible) would confirm or refute it directly,
    rather than leaving it as a plausible but unverified explanation.
  \item[Ground truth diversification] Using additional active-scanning ground truths beyond TLScope
    to validate the findings and reduce potential biases.
  \item[Indicator quality scoring] Implementing RAPIER-style \cite{qing2024rapier} confidence
    scoring to distinguish high-confidence from low-confidence provider detections.
\end{description}

% ============================================================
\chapter{Conclusion}
\label{ch:conclusion}
% ============================================================

This thesis presents NetScan, a longitudinal study of malicious IP indicator lifetimes across six
public threat intelligence providers. Using TLScope-classified malicious IP addresses as ground
truth, eight groups of 1\,596--2\,728 IPs each (18\,428 IPs in total) were tracked, on a
predominantly 3-day cadence, for one to two months each.

The study is motivated by a fundamental question: given that TLScope has already identified a
server as malicious, how does the public threat intelligence ecosystem respond? The three research
questions --- detection lag, indicator persistence, and inter-provider agreement --- together
characterise the temporal effectiveness of public feeds relative to a ground truth produced by
active TLS scanning.

The results demonstrate a substantial coverage gap: 81.2\% of TLScope-identified malicious IPs were
never flagged by any tracked provider. Among the IPs that were eventually detected, 44.0\% were
flagged with zero lag and 40.7\% only after missing the first collection cycle, and agreement
between detecting engines is low --- 83.9\% of ever-detected IPs are flagged by exactly one engine,
with the notable exception of BitDefender and G-Data, whose detections are identical
(Jaccard~=~1.0) in every single wave. Together these results give a quantitative, longitudinal
answer to where public threat intelligence feeds succeed and where they fall structurally short
relative to specialised active TLS scanning, and provide a concrete benchmark --- rather than a
prediction --- for future threat intelligence system evaluation.

Two further findings qualify the headline number. First, the never-detected rate is not flat over
the study period: it rises from $\sim$73\% in the earliest waves to a 82--87\% band from wave
4 onward, a trend that is consistent in direction but, across only eight waves, stops just short of
conventional significance (Pearson $r = 0.704$, $p = 0.051$) and is in any case confounded with
Censys being deactivated in the collection tool after wave 4 and never restored
(\S\ref{sec:provider_availability}) --- a methodological change this thesis discovered by auditing
the collected data rather than one that was flagged at the time. Second, among the five standalone
providers evaluated independently of VirusTotal's aggregated engines, Cinsscore detected none of the
18\,428 tracked IPs and OpenPhish detected only 2, so effectively all of this dataset's coverage of
TLScope-identified infrastructure traces back to VirusTotal and, to a lesser extent, ThreatFox and
AbuseIPDB (\S\ref{sec:provider_coverage}). Both findings point to the same conclusion as the
headline result: public feed coverage of active-scanning-identified malicious infrastructure is not
just incomplete but unevenly and inconsistently so, across providers and across time.

% ============================================================
%  BACK MATTER
% ============================================================

\appendix

\chapter{Configuration Reference}
\label{app:config}

API keys are stored in a \texttt{.env} file in the project root. Copy \texttt{.env\_ex} and fill
in the required credentials:
\begin{lstlisting}[language=bash, caption={Example .env configuration}]
VIRUSTOTAL = "VIRUSTOTAL Key::limit"
ABUSEIPDB  = "ABUSEIPDB Key::limit"
CENSYS     = "CENSYS Key||secret::limit"
THREATFOX  = "THREATFOX Key::limit"
# No API key needed for these services; set to "local"
CINSSCORE  = "local::10000"
OPENPHIS   = "local::10000"
\end{lstlisting}

Each value follows the form \texttt{key||secret::rate\_limit}. \texttt{runfetch.py} splits on
\texttt{::} to separate the credential from the per-service rate limit, then on \texttt{||} to
separate key from secret for the one provider that needs both (Censys). Cinsscore and OpenPhish
require no API key --- they use publicly accessible list endpoints --- but still need an entry, set
to \texttt{local}, because the service is registered by the presence of its environment variable.
Note that the OpenPhish variable is spelled \texttt{OPENPHIS} in the codebase; the name must match
exactly for the service to be picked up. Censys is present in \texttt{.env\_ex} but is not read in
the current configuration, since it is commented out of the active service list
(\S\ref{sec:provider_availability}).

\chapter{CLI Reference}
\label{app:cli}

\section{runfetch.py}

\begin{lstlisting}[language=bash]
python runfetch.py -f <folder> [-w WORKERS] [-vw VT_WORKERS]
                   [-q QUEUE_SIZE] [-sf START_FROM] [-dbg]
\end{lstlisting}

\texttt{-w} sets the number of producer workers (Section~\ref{ch:implementation}); \texttt{-vw}
sets the number of VirusTotal-polling consumer workers, independently of \texttt{-w}.

\section{runanalysis.py}

\begin{lstlisting}[language=bash]
python runanalysis.py -f <folder> [-t TOP_ENGINES] [-d DAY_DIFF] [-w WORKERS]
                      [-q QUEUE_SIZE] [-o OUTPUT_DIR] [--do_group_graphs]
\end{lstlisting}

\texttt{-t} sets how many of the highest peak-detection engines get per-engine figures and
summary rows (default 3; the analyses reported in Chapter~\ref{ch:results} were run with
\texttt{-t 5}, and the per-wave \texttt{top\_engines} value is recorded in each wave's
\texttt{out\_logs.txt}). Note that this setting affects only which engines are singled out for
figures and the \texttt{Report Stats} block: the aggregate, agreement and chance-baseline results
are computed over \emph{all} engines in \texttt{engine\_summary.csv} and are unchanged by it.
\texttt{--do\_group\_graphs}
additionally plots a per-IP malicious-vote figure for every group in the run; it is off by default
since it produces one figure per group rather than one aggregate figure per statistic, and is
impractical for combined multi-wave runs.

Configuration overrides are read at two levels. A \texttt{config.json} placed directly inside
\texttt{<folder>} overrides the corresponding command-line value or default for the whole run, for
any key matching a CLI argument name. A \texttt{config.json} placed inside an individual
\emph{sub-group} directory overrides that key for that sub-group only, leaving the rest of the wave
on the run-level value; unknown keys are reported and ignored. The per-sub-group level is the one
used in practice: it is how waves 1 and 5 were collected on 1-day and 2-day cadences
(\S\ref{sec:collection_schedule}) while the run default stayed at 3, and the resulting per-group
values are echoed into \texttt{out\_logs.txt} as \texttt{group\_overrides} so the effective cadence
of each sub-group is recoverable from the analysis output alone. Missing keys or a missing file
silently fall back to the CLI/default
value. This allows a group-specific setting (for example, an irregular actual scan cadence) to be
recorded once alongside the group's data rather than passed by hand on every invocation.

% --- Bibliography ---
\cleardoublepage
\addcontentsline{toc}{chapter}{Bibliography}

\begin{thebibliography}{99}

\bibitem{theofanous2024tlscope}
A.~Theofanous, E.~Papadogiannaki, A.~Shevtsov, and S.~Ioannidis,
``Fingerprinting the Shadows: Unmasking Malicious Servers with Machine Learning-Powered TLS Analysis,''
in \textit{Proc.\ ACM Web Conference (WWW)}, 2024, pp.\ 1933--1944.
DOI:~10.1145/3589334.3645719.

\bibitem{antonakakis2010notos}
M.~Antonakakis, R.~Perdisci, D.~Dagon, W.~Lee, and N.~Feamster,
``Building a Dynamic Reputation System for DNS,''
in \textit{Proc.\ 19th USENIX Security Symposium}, 2010, pp.\ 273--290.

\bibitem{bilge2011exposure}
L.~Bilge, E.~Kirda, C.~Kruegel, and M.~Balduzzi,
``EXPOSURE: Finding Malicious Domains Using Passive DNS Analysis,''
in \textit{Proc.\ NDSS 2011}, Feb.\ 2011.

\bibitem{deniz2025mantis}
F.~Deniz, M.~Nabeel, T.~Yu, and I.~Khalil,
``MANTIS: Detection of Zero-Day Malicious Domains Leveraging Low Reputed Hosting Infrastructure,''
in \textit{Proc.\ 2025 IEEE Symposium on Security and Privacy (S\&P)}, 2025, pp.\ 1789--1807.
DOI:~10.1109/SP61157.2025.00067.

\bibitem{li2025conceptdrift}
A.~S.~Li, A.~Iyengar, A.~Kundu, and E.~Bertino,
``Revisiting Concept Drift in Windows Malware Detection: Adaptation to Real Drifted Malware with Minimal Samples,''
in \textit{Proc.\ 32nd NDSS 2025}, Feb.\ 2025.

\bibitem{qing2024rapier}
Y.~Qing, Q.~Yin, X.~Deng, Y.~Chen, Z.~Liu, K.~Sun, K.~Xu, J.~Zhang, and Q.~Li,
``Low-Quality Training Data Only? A Robust Framework for Detecting Encrypted Malicious Network Traffic,''
in \textit{Proc.\ 31st NDSS 2024}, Feb.\ 2024.

\bibitem{tian2025density}
J.~Tian, W.~Kong, D.~Gao, T.~Wang, T.~Gu, K.~Qiu, Z.~Wang, and X.~Kuang,
``Density Boosts Everything: A One-stop Strategy for Improving Performance, Robustness, and Sustainability of Malware Detectors,''
in \textit{Proc.\ 32nd NDSS 2025}, Feb.\ 2025.

\bibitem{williams2024sixsense}
G.~Williams, M.~Erdemir, A.~Hsu, S.~Bhat, A.~Bhaskar, F.~Li, and P.~Pearce,
``6Sense: Internet-Wide IPv6 Scanning and its Security Applications,''
in \textit{Proc.\ 33rd USENIX Security Symposium}, 2024, pp.\ 2281--2298.

\bibitem{chen2024beam}
Y.~Chen, Q.~Yin, Q.~Li, Z.~Liu, K.~Xu, Y.~Xu, M.~Xu, Z.~Liu, and J.~Wu,
``Learning with Semantics: Towards a Semantics-Aware Routing Anomaly Detection System,''
in \textit{Proc.\ 33rd USENIX Security Symposium}, 2024, pp.\ 5143--5160.

\bibitem{ramesh2024calcuLatency}
R.~Ramesh, P.~Winter, S.~Korman, and R.~Ensafi,
``CalcuLatency: Leveraging Cross-Layer Network Latency Measurements to Detect Proxy-Enabled Abuse,''
in \textit{Proc.\ 33rd USENIX Security Symposium}, 2024, pp.\ 2263--2280.

\bibitem{zhang2024staleglue}
Y.~Zhang, B.~Liu, H.~Duan, M.~Zhang, X.~Li, F.~Shi, C.~Xu, and E.~Alowaisheq,
``Rethinking the Security Threats of Stale DNS Glue Records,''
in \textit{Proc.\ 33rd USENIX Security Symposium}, 2024, pp.\ 1261--1277.

\bibitem{zhang2024resolverfuzz}
Q.~Zhang, X.~Bai, X.~Li, H.~Duan, Q.~Li, and Z.~Li,
``ResolverFuzz: Automated Discovery of DNS Resolver Vulnerabilities with Query-Response Fuzzing,''
in \textit{Proc.\ 33rd USENIX Security Symposium}, 2024, pp.\ 4729--4746.

\bibitem{jarm2020}
Salesforce Engineering,
``JARM: A Solid TLS Fingerprinting Tool,''
\url{https://engineering.salesforce.com/easily-identify-malicious-servers-on-the-internet-with-jarm-e095edac525a}, 2020.

\end{thebibliography}

\end{document}
